\chapter{Averaging Protocol}

Consideriamo una rete modellata come un grafo connesso e non bipartito $G = (V, E)$. Ci troviamo sotto le restrizioni 
\[
    R = \{\texttt{TR, CN, BL, SYNC, LOCAL}\} 
\]
Ogni nodo $v \in V$ possiede un valore, diciamo $x(v) \in \mathbb{R}^+$ e localmente deve calcolare la media globale $\mu$ di tutti i valori posseduti da ciascun nodo
\[
    \mu=\frac{1}{|V|}\sum_{v\in V} x(v).
\]
Si richiede quindi una computazione distribuita della media: al termine, ogni nodo deve aver eseguito il calcolo di \(\mu\) a partire dalle sole informazioni scambiate nel protocollo.

Possiamo quindi modellare il problema con la tripla $\langle P_{init}, P_{final}, R\rangle$, dove

\begin{itemize}
    \item $P_{init}:$ la configurazione iniziale di ciascun nodo è $x^{(0)}(v)$. Se consideriamo tutti i valori inizali, allora possiamo definire un vettore $\mathbf{x} = \langle x^{(0)}(1), \dots, x^{(0)}(n)\rangle^T$
    \item $P_{final}$ la configurazione finale di ciascun nodo è $x^{(t)}:\ |x^{(t)} - x^{(t-1)}| \leq \epsilon$.
\end{itemize}

Definiamo quindi il protocollo \texttt{AVG} come segue:
\begin{enumerate}
    \item Al round $t = 0$, ogni nodo $v \in V$ possiede un valore iniziale $x^{(0)}(v) \in \mathbb{R}$
    \item Ad ogni round $t \geq 1$, ogni nodo $v \in V$ aggiora il suo valore $x^{(t)}(v)$ con la \textbf{media dei valori posseduti dai suoi vicini} al termine del round precedente $t - 1$, ossia
    \[
        x^{(t)}(v) = \frac{1}{\delta(v)} \sum_{u \in N(v)} x^{(t-1)}(u)
    \]
    \item Se $x^{(t)}:\ |x^{(t)} - x^{(t-1)}| \leq \epsilon$ allora \textbf{return} $x^{(t)}$ e \textbf{stop}.
\end{enumerate}

Definiamo $\forall v \in V$ il suo grado come $\delta(v) = |N(v)|$ (compreso se stesso). Diremo inoltre che il grafo è $\delta$\texttt{-regular} se $\delta(v) = \delta\ \forall v \in V$.

\section{Analisi del protocolo \texttt{AVG} per grafi $\delta$\texttt{-regular} e connessi}

\begin{teorema}[Message Complexity]
    La message complexity ad ogni round del protocollo \texttt{AVG} è $\theta(m)$.
    \begin{proof}
        Consideriamo il nodo $v \in V$ e contiamo il numero di messaggi che invia per round. Ad ogni rount $t,\ v$ invia $x^{(t)}(v)$ a ciascuno dei suoi $\delta(v)$ vicini, quindi in totale, abbiamo che il numero di messaggi scambiati ad ogni round è:
        \[
            M[\texttt{AVG}]_t = \sum_{v \in V} \delta(v) = \sum_{v \in V} \delta(v) = 2m \in \theta(m)
        \]
    \end{proof}
\end{teorema}

Vogliamo ora dimostrare che questo protocollo, converge ad un certo valore $\mu$ che è esattamente la media globale, e inoltre vogliamo determinare il tempo di convergenza.

\subsection{Convergenza}

\begin{teorema}[Convergenza]\label{teo:convergenza}
    Se il grafo $G = (V, E)$ è connesso, non bipartito e $\delta\texttt{-regular}$, allora il protocollo \texttt{AVG} converge in norma $L_2$ ad una configurazione stabile nella quale ogni nodo ha come valore 
    \[
        x(v) = \frac{1}{n}\sum_{v \in V} x^{(0)}(v)
    \]
    ossia, \textbf{converge} alla \textbf{media globale}.
\end{teorema}

La dimostrazione di questo teorema si articola in diverse fasi. In primo luogo, stabiliremo la convergenza nel caso più generale di un grafo non $\delta\texttt{-regular}$, per poi applicare l'ipotesi di regolarità.

Inoltre, per dimostrare questo teorema, faremo uso dell'algebra lineare e della teoria spettrale.

\begin{definizione}[Norma $L_p$]
    Sia $\mathbf{x} \in \mathbb{R}^n$, per ogni $p \in \mathbb{N}$ fissato, definiamo la norma $L_p$ come
    \[
        \norm{\mathbf{x}}_p = \left( \sum_{i = 1}^{n} |x_j|^p \right)^{1/p}
    \]
\end{definizione}

\begin{definizione}[Norma $L_\infty$]
    Sia $\mathbf{x} \in \mathbb{R}^n$, definiamo la norma $L_\infty$ come
    \[
        \norm{\mathbf{x}}_\infty = \max \{ |x_j|: j \in [n] \}
    \]
\end{definizione}

\begin{lemma}
    Per ogni vettore $\mathbf{x} \in \mathbb{R}^n$, vale che:
    \[
        \norm{\mathbf{x}}_\infty \leq \norm{\mathbf{x}}_2 \leq \norm{\mathbf{x}}_1 \leq \sqrt{n}\norm{\mathbf{x}}_2 \leq n\norm{\mathbf{x}}_\infty
    \]
\end{lemma}

\begin{definizione}[Base Ortonormale]
    Sia $V$ uno spazio vettoriale di dimensione finita sul campo $K$, nel quale sia definito il prodotto scalare. Una base ortonormale per $V$ è una base composta dai vettori $\langle \mathbf{u}_1, \dots, \mathbf{u}_n \rangle$ tale che:
    \begin{itemize}
        \item \textbf{Vettori mutualmente ortogonali:} $\langle \mathbf{u}_i, \mathbf{u}_j \rangle = 0, \quad \forall i \neq j$
        \item \textbf{Norma uguale a 1:} $\norm{\mathbf{u}_i}_2 = 1$
    \end{itemize}
\end{definizione}
Se $B = \langle \mathbf{u}_1, \dots, \mathbf{u}_n \rangle$ è ortonormale, allora
\[    
    \mathbf{x} = \sum_{\mathbf{u}_i \in B} \langle \mathbf{x}, \mathbf{u}_i \rangle \mathbf{u}_i
\]

Infine, definiamo la matrice di adiacenza $A,\ n \times n$ del grafo $G$ ponendo $A(i, j) = 1 \iff (i, j) \in E$. Definiamo inoltre la matrice diagonale $D^{-1},\ n \times n$ ponendo $D^{-1}(i, i) =  \frac{1}{\delta(i)}\ \forall i \in V$ e nelle altre entrate 0.

\begin{lemma}
    Per ogni round $t \geq 1$, la configurazione del sistema al round $t$ soddisfa la seguente equazione:
    \begin{equation}\label{eq:1}
        \mathbf{x}^{(0)} := \mathbf{x} \qquad \mathbf{x}^{(t+1)} := (D^{-1}A)\mathbf{x}^{(t)} = D^{-1} (A\mathbf{x}^t) 
    \end{equation}
    Più precisamente, se definiamo la matrice $P := D^{-1}A$ la \textbf{matrice di transizione} del grafo $G$, allora
    \begin{equation}\label{eq:2}
        \mathbf{x}^{(0)} := \mathbf{x} \qquad \mathbf{x}^{(t+1)} := P\mathbf{x}^{(t)} = PP\mathbf{x}^{(t-1)} = \dots = P^t\mathbf{x} \quad \text{Power Method}
    \end{equation}
    Questa equazione ci sta modellando un sistema dinamico lineare, in cui la configurazione al round $t + 1$ è ottenuta moltiplicando la matrice $P$ per la configurazione al round $t$.
    \begin{proof}
        Per ogni $i \in V$, possiamo scrivere la componente $i$-esima del vettore $\mathbf{x^{(t+1)}}$ come
        \[
            x^{(t+1)}(i) = \frac{1}{\delta(i)} \sum_{j = 1}^{n} A(i, j)x^{(t)}(j) = \frac{1}{\delta(i)} \sum_{j \in N(i)}x^{(t)}(j)
        \]
    \end{proof}
\end{lemma}

\textbf{Osservazione:} La matrice $P$ è \textbf{diagonalizzabile}, poichè è una matrice \textit{simile} ad una simmetrica.

\begin{definizione}[Matrici Simili]
    Due matrici $A, B,\ n \times n$ si dicono simili se esiste una matrice invertibile $M$ con la proprietà che $B = M^{-1}AM$
\end{definizione}

\begin{teorema}
    Se $A$ e $B$ sono matrici simili, allora:
    \begin{enumerate}
        \item $\det A = \det B$;
        \item $c_A(x) = c_B(x)$;
        \item $A$ e $B$ hanno gli stessi autovalori.
    \end{enumerate}
\end{teorema}

\begin{lemma}
    $P$ è una matrice simile alla matrice simmetrica $\tilde{P}$:
    \[
        \tilde{P} = D^{1/2}PD^{-1/2} = D^{1/2}\left(D^{-1}A\right)D^{-1/2} = D^{-1/2}AD^{-1/2}.
    \]
    Inoltre, $\tilde{P}$ è simmetrica, poiché:
    \[
        \tilde{P}^T = \left(D^{-1/2}AD^{-1/2}\right)^T = D^{-1/2}A^TD^{-1/2} = D^{-1/2}AD^{-1/2} = \tilde{P}.
    \]
\end{lemma}

Poichè $\tilde{P}$ è simmetrica, per il \textbf{Teorema Spettrale} ammette una base ortonormale di autovettori $\{\tilde{u}_1, \ldots, \tilde{u}_n\}$ con autovalori reali. Inoltre, poichè $P$ e $\tilde{P}$ sono simili, hanno gli stessi autovalori.

\textbf{Trasformazione degli autovettori.} Se $\tilde{u}$ è autovettore di $\tilde{P}$ con autovalore $\lambda$, ovvero $\tilde{P}\tilde{u} = \lambda\tilde{u}$, sostituendo la definizione di $\tilde{P}$:
\[
    D^{1/2}PD^{-1/2}\tilde{u} = \lambda\tilde{u}.
\]
Moltiplicando ambo i membri a sinistra per $D^{-1/2}$:
\[
    PD^{-1/2}\tilde{u} = \lambda D^{-1/2}\tilde{u}.
\]
Quindi $u = D^{-1/2}\tilde{u}$ è autovettore di $P$ con lo stesso autovalore $\lambda$. Poichè la trasformazione $\tilde{u} \mapsto D^{-1/2}\tilde{u}$ è invertibile (in quanto $D^{-1/2}$ è invertibile), gli $n$ autovettori $\{u_1, \ldots, u_n\}$ così ottenuti formano una \textbf{base di autovettori per} $P$, che risulta quindi diagonalizzabile.

\textbf{Osservazione:} La matrice $P$ è \textbf{stocastica}, ossia rispetta le seguenti proprietà:
\begin{itemize}
    \item \textbf{Non negatività degli elementi:}
    \[
        \forall i,j \in [n],\ P_{i,j} \geq 0
    \]
    \item \textbf{Somma unitaria per ogni riga:}
    \[
        \forall i \in [n],\ \sum_{j = 1}^{n} P_{ij} = 1
    \]
\end{itemize}
Inoltre, se $G$ è anche \textbf{regolare}, allora la matrice $P$ è \textbf{doppiamente stocastica}, $(P = P^T)$ ossia vale anche la \textbf{somma unitaria per ogni colonna:}
\[
    \forall j \in [n],\ \sum_{i = 1}^{n} P_{ij} = 1
\]

Applicando la teoria spettrale (decomposizione spettrale), possiamo individuare gli autovalori di $P$ e le loro proprietà.

Vogliamo studiare la convergenza del sistema in norma $L_2$.

\begin{lemma}
    Assumendo che $G$ è connesso e non bipartito e poichè la matrice $P$ è positiva a valori reali, ammette quindi $n$ autovalori $\lambda_1, \lambda_2, \dots, \lambda_n$ tali per cui $|\lambda_i| \leq 1,\ \forall i \in [n]$ e inoltre si ha un ordinamento di tali autovalori in senso non decrescente
    \[
        \langle \lambda_1 = 1, \lambda_2 < 1, \dots, \lambda_n \rangle
    \]
\end{lemma}

Ora, possiamo scrivere il vettore $\mathbf{x}$ come combinazione lineare della base formata dagli \textbf{autovettori} di $P$. Sia $\langle \mathbf{u}_1, \mathbf{u}_2, \dots, \mathbf{u}_n \rangle$ tale base, allora 
\[
    \mathbf{x} = \alpha_1\mathbf{u}_1 + \alpha_2\mathbf{u}_2 + \dots + \alpha_n\mathbf{u}_n
\]
e sostituendo nell'\autoref{eq:2}, otteniamo
\begin{equation}\label{eq:decomp-pt}
\begin{aligned}
    \mathbf{x}^{(t+1)}
    &= P^t \mathbf{x}
    = P^t\left(\alpha_1\mathbf{u}_1 + \alpha_2\mathbf{u}_2 + \dots + \alpha_n\mathbf{u}_n\right) \\
    &= \alpha_1 P^t\mathbf{u}_1 + \alpha_2 P^t\mathbf{u}_2 + \dots + \alpha_n P^t\mathbf{u}_n \\
    &= \alpha_1 \lambda_1^t \mathbf{u}_1 + \alpha_2 \lambda_2^t \mathbf{u}_2 + \dots + \alpha_n \lambda_n^t \mathbf{u}_n \\
    &= \alpha_1 \mathbf{u}_1 + \sum_{i=2}^{n} \alpha_i \lambda_i^t \mathbf{u}_i.
\end{aligned}
\end{equation}

\textbf{Osservazione:} Esistono autovettori con moltiplicità maggiore di 1, quindi è possibile che $\lambda_2 = \lambda_3 = \dots = \lambda_k$ per qualche $k > 2$. In questo caso, possiamo semplicemente raggruppare tali autovettori in un unico autovettore con autovalore $\lambda_2$.

Osserviamo quindi che per $t \to \infty$, siccome $|\lambda_i| < 1,\ \forall i > 1$, allora 
\[
    \sum_{i=2}^{n} \alpha_i \lambda_i^t \mathbf{u}_i \xrightarrow[t \to \infty]{} \mathbf{0}.
\]
in \textbf{norma}. 
Quindi, applicando un analisi spettrale della matrice di transizione $P$ del grafo, sotto certe condizioni abbiamo che il sistema converge ad uno stato stazionario $\alpha_1\mathbf{u}_1$.

Se sfruttiamo l'ipotesi di $\delta$\texttt{-regular} del grafo $G$ e quindi che la matrice $P$ è simmetrica, allora la base formata dagli autovettori $\langle \mathbf{u}_1, \mathbf{u}_2, \dots, \mathbf{u}_n \rangle$ è \textbf{ortonormale}, ossia i vettori $\mathbf{u}_i$ sono \textit{mutualmente} ortogonali avente norma $\norm{\mathbf{u}_i}_2 = 1$.

\begin{lemma}
    A partire da una qualsiasi configurazione iniziale $\mathbf{x}$, il sistema converge (in qualsiasi norma) a 
    \begin{equation}\label{eq:convergence}
        \mathbf{x}^{(t)} \to \alpha_1\mathbf{u}_1
    \end{equation}
\end{lemma}

Ci domandiamo ora chi sono $\alpha_1$ e $\mathbf{u}_1$. Per definizione di autovettore con autovalore $\lambda = 1$, abbiamo che $P \mathbf{u}_1 = \lambda \mathbf{u}_1 = \mathbf{u}_1$,
\[
\mathbf{u}_1 = P\,\mathbf{u}_1 =
\begin{bmatrix}
\frac{1}{\delta_1} & 0 & \cdots & \frac{1}{\delta_1} & 0 & \cdots & \frac{1}{\delta_1} \\
0 & \frac{1}{\delta_2} & \cdots & \frac{1}{\delta_2} & 0 & \cdots & 0 \\
\vdots & \vdots & \ddots & \vdots & \vdots & \ddots & \vdots \\
0 & 0 & \cdots & \frac{1}{\delta_i} & \cdots & \frac{1}{\delta_i} & 0 \\
0 & 0 & \cdots & \frac{1}{\delta_n} & \cdots & 0 & \frac{1}{\delta_n}
\end{bmatrix}
\begin{bmatrix}
u_1(1) \\
u_1(2) \\
\vdots \\
u_1(i) \\
\vdots \\
u_1(n)
\end{bmatrix}.
\]

Poiché a noi ci interessa la massa di $\mathbf{x}$ (ossia $\sum_{j \in V} x(j)$), possiamo considerare un qualsiasi autovettore $\mathbf{v}_1$ relativo a $\lambda = 1$ e lo possiamo normalizzare (ottenere un nuovo autovettore di norma 1) $\mathbf{u}_1$ nel seguente modo:
\[
    \mathbf{u}_1 = \frac{\mathbf{v_1}}{\norm{\mathbf{v}_1}_2}
\]
In questo modo, otteniamo che $\alpha_1 = \langle \mathbf{x}, \mathbf{u}_1 \rangle$.

Sostituendo ora nell'\autoref{eq:convergence}, otteniamo che
\begin{equation}\label{eq:convergence2}
    \mathbf{x}^{(t)} \to \alpha_1\mathbf{u}_1 = \langle \mathbf{x}, \mathbf{u_1} \rangle \mathbf{u}_1 = \frac{\mathbf{v_1}}{\norm{\mathbf{v}_1}_2} \cdot \left( \sum_{j = 1}^{n} x(j)u_1(j) \right)
\end{equation}
In questo modo, ci rimane solo da determinare chi è il nostro autovettore $\mathbf{v}_1$ relativo a $\lambda = 1$.

\textbf{Osservazione:} $\alpha_i = \langle \mathbf{x}, \mathbf{u}_i \rangle$ sono le proiezioni del vettore iniziale $\mathbf{x}$ lungo la base (\textit{ortonormale}) di vettori. Inolre, $|\alpha_i| = |\langle \mathbf{x}, \mathbf{u}_i \rangle| \leq \mathcal{M} = \sum_{j \in V} x(j)$.

Ora ci chiediamo chi è $\mathbf{v}_1$?

Abbiamo due casistiche, caso 1: $P$ è simmetrica, caso 2: $P$ non è simmetrica.

\textbf{CASO 1: $\mathbf{G}$ non regolare, quindi $\mathbf{P}$ non simmetrica.}

\begin{lemma}\label{lemma:right_eigenvector}
    Il vettore $\mathbf{y} = \langle \frac{\delta_1}{2m}, \frac{\delta_2}{2m}, \dots, \frac{\delta_n}{2m} \rangle$ è un \textbf{1-autovettore destro} della matrice $P^T = (D^{-1}A)^T = AD^{-1}$, ossia
    \[
        P^T\mathbf{y} = \mathbf{y}
    \]
    \begin{proof}
        Calcoliamo l'$i$-esimo elemento del prodotto $P^T\mathbf{y}$
        \begin{align*}
            \left(P^T\mathbf{y} \right)_i &= \sum_{j = 1}^{n} \left(P^T\right)_{i,j}\mathbf{y}_j
            = \sum_{j = 1}^{n} \frac{A(i, j)}{\delta_j}\frac{\delta_j}{2m} \\
            &= \sum_{j = 1}^{n} \frac{A(i, j)}{2m} = \frac{1}{2m}\sum_{j = 1}^{n} A(i, j) \\
            &= \frac{\delta_i}{2m}
        \end{align*}
        Poiché $\sum_{j = 1}^{n} A(i, j)$ rappresenta la somma della riga $i$ della matrice di adiacenza, ovvero il numero di totale di collegamenti del nodo $i$, ce per definizione è il suo grado $\delta_i$.
    \end{proof}
\end{lemma}
In sintesi, abbiamo individuato un autovettore \textbf{destro} di $P^T$, cioè $\mathbf{y}$ tale che $P^T\mathbf{y}=\mathbf{y}$. Inoltre, dalla relazione
\[
\left(P^T\mathbf{y}\right)^T=\mathbf{y}^TP,
\]
segue che $\mathbf{y}^T$ è un autovettore \textbf{sinistro} di $P$ associato all'autovalore $1$.

Tuttavia, per la nostra analisi serve un autovettore \textbf{destro} di $P$. In generale, se $P\neq P^T$, autovettori sinistri e destri non coincidono. Possiamo però ottenerne uno in modo esplicito ponendo
\[
\mathbf{v}_1:=D^{-1}\mathbf{y}.
\]
Mostriamo ora che $\mathbf{v}_1$ è effettivamente un autovettore destro di $P$ relativo a $\lambda=1$:

\[
    P\mathbf{v}_1 = (D^{-1}A)D^{-1}\mathbf{y} = D^{-1}(AD^{-1})\mathbf{y} = D^{-1}P^T\mathbf{y} = D^{-1}\mathbf{y} = \mathbf{v}_1
\]
dove
\[
    \mathbf{v}_1 = D^{-1}\langle \frac{\delta_1}{2m}, \frac{\delta_2}{2m}, \dots, \frac{\delta_n}{2m} \rangle = \langle \frac{1}{2m}, \frac{1}{2m}, \dots, \frac{1}{2m} \rangle
\]
Abbiamo quindi individuato un autovettore \textbf{destro} $\mathbf{v}_1$, valido anche nel caso in cui il grafo $G$ non sia regolare. Tuttavia, nel nostro caso stiamo considerando grafi $\delta$\texttt{-regular}, cioè grafi in cui ogni nodo ha lo stesso grado $\delta_i = \delta$. 

\textbf{CASO 2: $\mathbf{G}$ regolare, quindi $\mathbf{P}$ simmetrica.}

\begin{lemma}
    Se $G$ è $\delta$\texttt{-regular} (quindi $P$ è simmetrica), allora un autovettore destro di $P$ associato a $\lambda=1$ è
    \[
    \mathbf{v}_1=\left\langle \frac{\delta}{2m},\frac{\delta}{2m},\dots,\frac{\delta}{2m}\right\rangle
    \]

    \begin{proof}
    Dal \autoref{lemma:right_eigenvector} sappiamo che
    \[
    \mathbf{y}=\left\langle \frac{\delta_1}{2m},\frac{\delta_2}{2m},\dots,\frac{\delta_n}{2m}\right\rangle
    \]
    soddisfa
    \[
    P^T\mathbf{y}=\mathbf{y}.
    \]
    Se il grafo è $\delta$\texttt{-regular}, allora $\delta_i=\delta$ per ogni $i$, quindi
    \[
    \mathbf{y}=\left\langle \frac{\delta}{2m},\frac{\delta}{2m},\dots,\frac{\delta}{2m}\right\rangle.
    \]
    Inoltre, nel caso regolare $P$ è simmetrica, dunque $P=P^T$. Ne segue
    \[
    P\mathbf{y}=P^T\mathbf{y}=\mathbf{y}.
    \]
    Pertanto $\mathbf{y}$ (cioè $\mathbf{v}_1$) è un autovettore destro di $P$ relativo all'autovalore $1$.
    \end{proof}
\end{lemma}

Ne segue che possiamo dimostrare il \autoref{teo:convergenza} come segue:
\begin{proof}
    Sappiamo che il sistema converge a $\alpha_1\mathbf{u}_1$. Essendo una base ortonormale, $\alpha_1 = \langle \mathbf{x}, \mathbf{u_1} \rangle$, quindi
    \[
        \mathbf{x}^{(t)} \to \langle \mathbf{x}, \mathbf{u}_1 \rangle \mathbf{u}_1
    \]
    Ora, avevamo definito $\mathbf{u}_1 = \frac{\mathbf{v}_1}{\norm{\mathbf{v_1}}_2}$, sostituendo 2 volte questa definizione otteniamo
    \[
        \mathbf{x}^{(t)} \to \left\langle \mathbf{x}, \frac{\mathbf{v}_1}{\norm{\mathbf{v_1}}_2} \right\rangle \frac{\mathbf{v}_1}{\norm{\mathbf{v_1}}_2}
    \]
    Siccome la norma $\norm{\mathbf{v_1}}_2$ è uno scalare, per le proprietà del prodotto scalare la possiamo tirare fuori, ottenendo il quadrato della norma
    \[
        \mathbf{x}^{(t)} \to \frac{1}{\norm{\mathbf{v_1}}_2^2}\langle \mathbf{x}, \mathbf{v}_1 \rangle \mathbf{v}_1
    \]
    \begin{itemize}
        \item \textbf{Denominatore:} Per calcolare $\norm{\mathbf{v_1}}_2^2$ dobbiamo elevare al quadrato ogni componente di $\mathbf{v}_1$ e sommarle. Siccome ci sono $n$ componenti tutte uguali a $\frac{\delta}{2m}$, ottenimo:
        \[
            \norm{\mathbf{v_1}}_2^2 = n \left( \frac{\delta}{2m} \right)^2 = \frac{n\delta^2}{4m^2}
        \]
        \item \textbf{Prodotto Scalare} $\langle \mathbf{x}, \mathbf{v}_1 \rangle$:
        \[
            \langle \mathbf{x}, \mathbf{v}_1 \rangle = \sum_{j = 1}^{n}x(j)v_1(j) = \sum_{j = 1}^{n}x(j)\frac{\delta}{2m}
        \]
        \item Volendo calcolare il valore di un singolo nodo $i$, prendiamo l'$i$-esima componente del vettore $\mathbf{v}_1(i) = \frac{\delta}{2m}$
    \end{itemize}
    Mettendo tutto insieme per il nodo $i$, otteniamo:
    \begin{align*}
        x^{(t)}(i) &\to \frac{1}{\frac{n\delta^2}{4m^2}} \left( \sum_{j = 1}^{n}x(j)\frac{\delta}{2m} \right) \frac{\delta}{2m} \\
        &= \frac{1}{\frac{n\delta^2}{4m^2}} \left( \frac{\delta}{2m} \frac{\delta}{2m} \right) \sum_{j = 1}^{n}x(j) \\
        &= \frac{1}{\frac{n\delta^2}{4m^2}} \left( \frac{\delta^2}{4m^2} \right) \sum_{j = 1}^{n}x(j) \\
        &= \frac{4m^2}{n\delta^2} \left( \frac{\delta^2}{4m^2} \right) \sum_{j = 1}^{n}x(j) \\
        &= \frac{1}{n}\sum_{j = 1}^{n}x(j)
    \end{align*}
\end{proof}

\newpage

\subsection{Tempo di Convergenza}
\begin{teorema}\label{teo:terminazione}
    Il protocollo \texttt{AVG} termina dopo
    \[
    \mathcal{T} \;=\; \log\left(\frac{n\mathcal{M}}{\epsilon}\right) / \log\left(\frac{1}{\lambda}\right)
    \]
    iterazioni; inoltre, per ogni nodo \(i \in V\), al tempo \(\mathcal{T}\) vale
    \[
    \left|x^{(\mathcal{T})}(i)-x^{(\mathcal{T}-1)}(i)\right| \le \epsilon .
    \]
    \begin{proof}
        Poiché consideriamo il valore \textit{locale} di ciacun nodo, ci interessa utilizzare la $\norm{\cdot}_{\infty}$ in quanto ci dà un upper bound \textit{massimale} al valore dei nodi. Per ogni $t \leq 1$, abbiamo che
        \begin{equation}\label{eq:tempo_convergenza}
        \begin{aligned}
            \mathbf{x}^{(t+1)} - \mathbf{x}^{(t)} &= P^{t+1}\mathbf{x} - P^t\mathbf{x} \\
            &= P^{t+1}(\alpha_1\mathbf{u}_1 + \dots + \alpha_n\mathbf{u}_n) + P^{t}(\alpha_1\mathbf{u}_1 + \dots + \alpha_n\mathbf{u}_n) \\
            [P\mathbf{u}_i = \lambda_1\mathbf{u}_i] &= \lambda_1^{t+1}\alpha_1\mathbf{u}_1 + \dots + \lambda_1^{t+1}\alpha_n\mathbf{u}_n + \lambda_1^{t}\alpha_1\mathbf{u}_1 + \dots + \lambda_1^{t}\alpha_n\mathbf{u}_n \\
            &= \alpha_1\mathbf{u}_1 - \alpha_1\mathbf{u}_1 + \sum_{i = 2}^{n} (\lambda_i - 1)\lambda_i^t\alpha_i\mathbf{u}_i \\
            &= (\lambda_2 - 1)\lambda_2^t\alpha_2\mathbf{u}_2 + \dots + (\lambda_n - 1)\lambda_n^t\alpha_n\mathbf{u}_n
        \end{aligned}
        \end{equation}

        Definiamo ora 
        \[
            \lambda = \max \{ |\lambda_j|:\ j \geq 2 \}, \qquad \lambda_{\min} = \min \{ |\lambda_j|:\ j \geq 2 \}
        \]
        e ricordando che $G$ è connesso e non bipartito abbiamo che $0 < \lambda_{\min} \leq \lambda < 1$. Applicando quindi la $\norm{\cdot}_{\infty}$ all'equazione \autoref{eq:tempo_convergenza}, otteniamo
        \begin{align*}
            \norm{\mathbf{x}^{(t+1)} - \mathbf{x}^{(t)}}_{\infty} 
            &= \norm{P^{t+1}\mathbf{x} - P^t\mathbf{x}}_{\infty} \\
            &\leq n(1 - \lambda_{\min})\lambda^t\alpha_n\norm{\mathbf{u}_n}_{\infty} \\
            &= n(1 - \lambda_{\min})\lambda^t\mathcal{M}
        \end{align*}
        Applicando la norma, $\norm{\lambda_{\min} - 1}_{\infty} = |\lambda_{\min} - 1| = 1 - \lambda_{\min}$. Ricordiamo che $\mathcal{M} = \sum_{j \in V} x(j)$ e osserviamo che $\forall i \in [n],\ \alpha_i \leq \mathcal{M}$. Quindi possiamo individuare quanto debba essere $\mathcal{T}$ affichè l'ultimo incremento sia sufficientemente piccolo, ossia
        \begin{align*}
            \norm{\mathbf{x}^{(t+1)} - \mathbf{x}^{(t)}}_{\infty} 
            &\leq n(1 - \lambda_{\min})\lambda^\mathcal{T}\mathcal{M} \leq \epsilon \\
            &\implies \lambda^\mathcal{T} \leq \frac{\epsilon}{n\mathcal{M}(1 - \lambda_{\min})} \\
            &\implies \mathcal{T} \log(\lambda) \leq \log\left(\frac{\epsilon}{n\mathcal{M}(1 - \lambda_{\min})}\right) \\
            &\implies \mathcal{T} \geq \frac{\log\left(\frac{\epsilon}{n\mathcal{M}(1 - \lambda_{\min})}\right)}{\log(\lambda)} \quad \text{(poiché } |\lambda| < 1 \implies \log(\lambda) < 0 \text{)} \\
            &\implies \mathcal{T} \geq \frac{-\log\left(\frac{n\mathcal{M}(1 - \lambda_{\min})}{\epsilon}\right)}{-\log\left(\frac{1}{\lambda}\right)} \\
            &\implies \mathcal{T} \geq \frac{\log\left(\frac{n\mathcal{M}(1 - \lambda_{\min})}{\epsilon}\right)}{\log\left(\frac{1}{\lambda}\right)} \\
            &\implies \mathcal{T} \geq \frac{\log\left(\frac{n\mathcal{M}}{\epsilon}\right)}{\log\left(\frac{1}{\lambda}\right)} \quad \text{(poiché } |\lambda_{\min}| < 1 \text{)} \\ 
        \end{align*}
    \end{proof}
\end{teorema}

\textbf{Osservazione:} Il tempo di convergenza è logaritmico rispetto a $n$ e $\mathcal{M}$, quindi nella dimensione del grafo e nella massa totale (\textbf{molto buono!}) e inversamente proporzionale a $\log(1/\lambda)$, ossia dipende dalla distanza tra 1 e il secondo autovalore $\lambda$ cioè dalla \textbf{spectral gap} (topologia) del grafo.
\begin{itemize}
    \item Nel caso di grafi regolari, determina la qualità d'espansione del grafo, ossia quanto è "ben connesso" il grafo. Più il grafo è un buon expander, più grande è la spectral gap ($\lambda_1 - \lambda_2$ è grande), quindi il grafo è ben connesso e il protocollo converge più velocemente. Se il grafo è un buon expander, allora $\lambda$ è limitato superiormente da una costante strettamente minore di 1, indipendentemente dal numero di nodi $n$, garantendo una convergenza rapida.
    \item Nel caso di grafi non regolari, la spectral gap è influenzata anche dalla distribuzione dei gradi dei nodi, quindi se ci sono nodi con gradi molto diversi tra loro, la spectral gap potrebbe essere più piccola, rallentando la convergenza del protocollo.
\end{itemize}

\begin{teorema}[Terminazione e Corretezza (Accuratezza)]
    Dato in input un qualsiasi grafo $G = (V, E)$ connesso, non bipartito e $\delta$\texttt{-regular} e un parametro di confidenza $\epsilon > 0$, il protocollo \texttt{AVG} termina dopo
    \[
        \mathcal{T} \;=\; \log\left(\frac{n\mathcal{M}}{\epsilon}\right) / \log\left(\frac{1}{\lambda}\right)
    \]
    passi. Inoltre, ogni nodo $i \in [n]$ restituisce un valore $x^{\mathcal{T}}(i)$ tale che:
    \[
        \left| x^{\mathcal{T}}(i) - \frac{1}{n}\sum_{i \in V} x(i) \right| \leq \epsilon
    \]
    \begin{proof}
        Come abbiamo visto sopra, possiamo limitare la variazione tra due passi successivi. Sia $t \geq 1$, allora vale sempre che
        \[
            \norm{\mathbf{x}^{(t+1)} - \mathbf{x}^{(t)}}_{\infty} \leq n\lambda^t\mathcal{M}
        \]
        Limitiamo poi la distanza della vera media globale. Andiamo a guardare la variazione tra il valore attuale $\mathbf{x}^{(t)}$ e il limite reale $\alpha_1\mathbf{u}_1$.
        \begin{align*}
            \norm{\mathbf{x}^{(t)} - \alpha_1\mathbf{u}_1}_{\infty} &= \norm{P^{t}\mathbf{x} - \alpha_1\mathbf{u}_1}_{\infty} \\
            &= \lambda_1^{t}\alpha_1\mathbf{u}_1 + \dots + \lambda_n^{t}\alpha_n\mathbf{u}_n - \alpha_1\mathbf{u}_1 \\
            [\lambda_1 = 1] &= (\alpha_1\mathbf{u}_1 - \alpha_1\mathbf{u}_1) + \lambda_2^{t}\alpha_2\mathbf{u}_2 + \dots + \lambda_n^{t}\alpha_n\mathbf{u}_n \\
            &=  \lambda_2^{t}\alpha_2\mathbf{u}_2 + \dots + \lambda_n^{t}\alpha_n\mathbf{u}_n \\
            &\leq n\lambda^t\mathcal{M}\norm{\mathbf{u}_n}_\infty = n\lambda^t\mathcal{M}
        \end{align*}

        Grazie al \autoref{teo:terminazione}, possiamo affermare che, quando il nodo $i$ termina la sua computazione, dopo $\mathcal{T}$ round, il suo valore è:
        \[
            \left|x^{(\mathcal{T})}(i) - \alpha_1\mathbf{u}_1\right| = \left|x^{(\mathcal{T})}(i) - \frac{\mathcal{M}}{n}\right| \leq n\lambda^{\mathcal{T}}\mathcal{M} \leq \epsilon
        \]
    \end{proof}
\end{teorema}